\documentclass[11pt]{letter}
\usepackage[utf8]{inputenc}
\usepackage[T1]{fontenc}
\usepackage[english]{babel}
\usepackage[a4paper,margin=1in]{geometry}

\signature{Bartosz Lewandowski, Ph.D.\\
Department of Clinical Engineering, Academy of Silesia in Katowice\\
Rolna 43, 40-555 Katowice, Poland\\
bartosz.lewandowski@akademiaslaska.pl}
\address{}

\begin{document}
\begin{letter}{The Editor-in-Chief\\ \textit{Computers in Human Behavior}}

\opening{Dear Editor,}

We submit the manuscript entitled ``Employee response accuracy on repeated questions: forgetting and
persistent errors in a three-year microlearning program'' for consideration as a research article in
\textit{Computers in Human Behavior}.

The study examines how people learn and forget when a workplace training application repeatedly
presents the same multiple-choice questions. Research on the testing and spacing effects shows that
retrieval practice distributed over time improves retention, and laboratory studies show that errors
can persist despite feedback. Evidence on these processes in computer-delivered workplace training,
however, comes mainly from short studies with small samples. We analyzed 57,997 responses that
180 bank employees gave to 164 questions over three years. Because the application repeated questions
on a schedule that did not depend on the previous response, error correction and forgetting could be
compared without the selection introduced by adaptive algorithms.

The main findings are as follows. Accuracy increased with each repetition of a question, including in
a cohort of the same employee--question sequences, because error corrections outnumbered forgetting
events. Forgetting increased with the interval between repetitions and decreased with the number of
consecutive correct responses. Corrected errors were the least durable, and when an error recurred,
employees usually chose the same distractor as before, more often than expected from distractor
popularity. This pattern is consistent with persistent misconceptions rather than guessing.

We believe the manuscript fits the scope of the journal because it concerns the psychological
processes of learning, forgetting and error persistence in people using a computer-based training
system, and it draws practical conclusions for the design of such systems: a schedule adapted to the
learner's response history and feedback that addresses the selected distractor.

We confirm that this work has not been published previously, is not under consideration for
publication elsewhere, and that its publication has been approved by all authors. One author is an
employee of the bank in which the study was conducted; this is declared in the manuscript files. The
data are the property of the bank and cannot be shared publicly.

\closing{Yours sincerely,}

\end{letter}
\end{document}
